\documentclass[12pt]{article}
% Préambule commun à tous les documents extraits de « La Revue CAPES »
\usepackage[T1]{fontenc}
\usepackage[utf8]{inputenc}
\usepackage{lmodern}
\usepackage{amsmath,amssymb,amsthm,amsfonts,mathrsfs}
\usepackage{setspace}
\usepackage{listings}
\usepackage{pgfplots}
\pgfplotsset{compat=1.18}
\usepackage{longtable}
\usepackage{adjustbox}
\usepackage{lettrine}
\usepackage{tkz-tab}
\usepackage{changepage}
\usepackage{pdflscape}
\usepackage{graphicx}
\usepackage{tikz}
\usepackage{xcolor}
\usepackage[a4paper,top=2cm,bottom=2.5cm,left=2.5cm,right=2cm]{geometry}
\usepackage{fancyhdr}
\usepackage{draftwatermark}
\SetWatermarkText{La RevueCapes}
\SetWatermarkLightness{0.9}
\SetWatermarkScale{2}
\usepackage{hyperref}
\hypersetup{colorlinks=true,linkcolor=blue!50!black,urlcolor=blue!50!black}

\definecolor{revue}{RGB}{20,60,120}

\pagestyle{fancy}
\fancyhf{}
\renewcommand{\headrulewidth}{0.4pt}
\setlength{\headheight}{15pt}

% \entete{Matière}{Titre}{Type de document}
\newcommand{\entete}[3]{%
  \fancyhead[L]{\small\textsc{La Revue CAPES} -- #1}%
  \fancyhead[R]{\small #2}%
  \fancyfoot[C]{\thepage}%
  \thispagestyle{plain}%
  \begin{center}
    {\color{revue}\rule{\textwidth}{1.2pt}}\\[4pt]
    {\small\textsc{Concours directs CAP-CEG / CAPES -- Burkina Faso}}\\[6pt]
    {\LARGE\bfseries\color{revue} #2}\\[6pt]
    {\large\itshape #1 -- #3}\\[2pt]
    {\color{revue}\rule{\textwidth}{1.2pt}}
  \end{center}
  \vspace{0.5cm}%
}

% Encadré pour les remarques ajoutées dans les corrigés
\newcommand{\remarque}[1]{\par\medskip\noindent\fcolorbox{revue}{revue!6}{\parbox{\dimexpr\linewidth-2\fboxsep-2\fboxrule}{\textbf{\color{revue}Remarque.} #1}}\par\medskip}


\begin{document}
\entete{Culture générale}{Corrigé du sujet type de dissertation n°5}{Correction}
\begin{spacing}{1.5}

\textbf{Introduction}

Le coup d'État de 2022 au Burkina Faso, qui a vu le capitaine Ibrahim Traoré accéder au pouvoir, a suscité des débats intenses sur sa nature et ses implications. Certains y voient une rupture avec les principes démocratiques établis, tandis que d'autres le perçoivent comme une réponse pragmatique à une situation sécuritaire et socio-politique critique. Cette analyse vise à examiner les deux perspectives pour mieux comprendre les enjeux de cet événement majeur.

1. Le coup d'État : une rupture démocratique

1.1. L'interruption du processus démocratique

Le renversement du pouvoir civil issu d'élections démocratiques constitue une violation des normes constitutionnelles.

Exemple : Le président Roch Marc Christian Kaboré, bien que critiqué pour son incapacité à gérer l'insécurité, avait été élu par le peuple dans un cadre démocratique.
Analyse : Ce coup d'État fragilise davantage les institutions démocratiques dans un pays où la gouvernance civile reste fragile.

1.2. Un recul pour la stabilité institutionnelle

Les transitions militaires entraînent souvent une instabilité politique et institutionnelle.

Exemple : Depuis 2015, le Burkina Faso a connu plusieurs changements de régime non démocratiques, sapant la confiance des citoyens envers leurs institutions.
Analyse : Cette instabilité complique la mise en place de réformes durables nécessaires à la bonne gouvernance.

1.3. Les risques de légitimité à l'international

Les partenaires internationaux, souvent attachés aux principes démocratiques, ont vu ce coup comme une rupture inacceptable.

Exemple : La suspension du Burkina Faso par certaines organisations comme la CEDEAO montre que ce coup d'État a terni l'image du pays sur la scène internationale.

Analyse : Cette situation peut limiter l'accès à des financements et à une coopération cruciale pour la lutte contre l'insécurité.

2. Le coup d'État : une réponse nécessaire à l'effondrement sécuritaire

2.1. Une incapacité criante de l'ancien régime à gérer l'insécurité
Avant le coup d'État, le Burkina Faso faisait face à une crise sécuritaire grandissante, marquée par des attaques terroristes incessantes et l'effondrement des services publics.

Exemple : En 2022, près de 40/100 du territoire était hors du contrôle de l'État, avec des millions de déplacés internes.

Analyse : Le capitaine Ibrahim Traoré a justifié son action comme une nécessité pour sauver le pays d'un effondrement total.

2.2. Un renforcement immédiat des efforts militaires

La transition dirigée par Traoré a mis un accent particulier sur le redéploiement des forces armées et l'implication des volontaires pour la défense de la patrie (VDP).

Exemple : Sous sa direction, plusieurs initiatives ont été lancées pour regagner les zones occupées par les groupes terroristes.

Analyse : Cette approche a rassuré une partie de la population qui se sentait abandonnée sous le régime précédent.

2.3. Une réponse populaire et nationale.

L'arrivée d’Ibrahim Traoré au pouvoir a été perçue par de nombreux Burkinabè comme un acte patriotique.

Exemple : Les manifestations de soutien dans plusieurs régions montrent que ce coup d'État a bénéficié d'une certaine adhésion populaire.

Analyse : Le sentiment d'urgence sécuritaire semble avoir primé sur les principes démocratiques pour une grande partie de la population.

3. Les dilemmes entre rupture démocratique et réponse nécessaire

3.1. La tension entre urgence et légitimité

Bien que justifié par l'urgence sécuritaire, le coup d'État soulève des questions sur la durabilité de solutions imposées par la force.

Exemple : Si la sécurité s'améliore sous Traoré, il faudra tout de même préparer un retour à un ordre constitutionnel pour éviter un cycle de coups d'État.

3.2. La nécessaire réconciliation entre démocratie et sécurité
La gestion de la transition doit concilier les impératifs sécuritaires avec un plan clair pour rétablir la démocratie.

Exemple : Les promesses d'organiser des élections d'ici 2025 seront un test pour évaluer la volonté de la transition à restaurer l'ordre démocratique.

 \textbf{Conclusion}
 
Le coup d'État de 2022 au Burkina Faso dirigé par le capitaine Ibrahim Traoré peut être perçu à la fois comme une rupture démocratique et une réponse nécessaire à l'effondrement sécuritaire. Si la transition a permis de redonner un souffle nouveau dans la lutte contre le terrorisme, elle soulève des interrogations sur l'avenir politique du pays. La clé réside dans la capacité du régime à équilibrer les exigences sécuritaires et la restauration de l'ordre démocratique, en évitant une dépendance excessive à des solutions militaires.



\quad
\quad

\end{spacing}
\end{document}
